\uuid{N9p6}
\titre{Image réciproque}
\auteur{bodin ; hanani ; mzari}
\organisation{exo7}
\licence{CC BY-NC-SA 4.0}
\datecreate{2020-01-20}
\niveau{L1}
\module{Algèbre}
\chapitre{Logique, ensemble, raisonnement}
\sousChapitre{Ensemble}
\difficulte{3}
\theme{Ensembles, applications, image réciproque}
\exo7qcmid{lille-1.ensembles.024}

\contenu{
\begin{question}
\qclassification{100.02,101.01,102.01,102.02}
\qtags{motcle=image réciproque}
Soient $E$ et $F$ deux ensembles non vides et $f$ une application de $E$ dans $F$. Soit $B$ un sous-ensemble de $F$. Quelles sont les bonnes réponses ?
\begin{answers}
    \bad{$B=f(f^{-1}(B))$}
    \bad{$B\subset f(f^{-1}(B))$}
    \good{$f(f^{-1}(B))\subset B$}
    \bad{$f(f^{-1}(B))=F\setminus B$}
\end{answers}
\begin{explanations}
Soit $y\in f(f^{-1}(B))$. Donc il existe $x\in f^{-1}(B)$ tel que $y=f(x)$. Mais, $x\in f^{-1}(B)\Leftrightarrow f(x)\in B$. Donc $y=f(x)\in B$.
\end{explanations}
\end{question}
}
